
\documentclass[12pt]{article}

\usepackage{amsmath, amsthm, amssymb, mathtools}
\usepackage{geometry}
\usepackage{hyperref}
\usepackage{enumitem}
\usepackage{microtype}
\usepackage{xcolor}

\geometry{margin=1.25in}

\newtheorem{theorem}{Theorem}[section]
\newtheorem{proposition}[theorem]{Proposition}
\newtheorem{lemma}[theorem]{Lemma}
\newtheorem{corollary}[theorem]{Corollary}
\newtheorem{conjecture}[theorem]{Conjecture}
\theoremstyle{definition}
\newtheorem{definition}[theorem]{Definition}
\newtheorem{example}[theorem]{Example}
\theoremstyle{remark}
\newtheorem{remark}[theorem]{Remark}
\newtheorem{warning}[theorem]{Warning}

\DeclareMathOperator{\supp}{supp}
\DeclareMathOperator{\type}{type}
\DeclareMathOperator{\Stab}{Stab}
\newcommand{\R}{\mathbb{R}}
\newcommand{\C}{\mathbb{C}}
\newcommand{\Z}{\mathbb{Z}}
\newcommand{\N}{\mathbb{N}}
\newcommand{\T}{\mathcal{T}}
\newcommand{\PW}{\mathrm{PW}}
\newcommand{\Sph}{\mathbb{S}}
\newcommand{\wh}{\widehat}
\newcommand{\abs}[1]{\left|#1\right|}
\newcommand{\norm}[1]{\left\|#1\right\|}
\newcommand{\ip}[2]{\left\langle #1,\, #2\right\rangle}
\newcommand{\re}{\operatorname{Re}}

% Status markers
\newcommand{\proven}{\textbf{\color{teal}[Proved]}}
\newcommand{\conditional}{\textbf{\color{orange}[Conditional]}}
\newcommand{\open}{\textbf{\color{red}[Open]}}
\newcommand{\heuristic}{\textbf{\color{purple}[Heuristic]}}
\newcommand{\review}[1]{\par\smallskip\noindent\textbf{\color{red}[Review comment.]} {\color{red}#1}\par\smallskip}
\newcommand{\response}[1]{\par\noindent\textbf{\color{blue}[Response.]} {\color{blue}#1}\par\smallskip}
\newcommand{\resolved}{\textbf{\color{teal}[RESOLVED]}}
\newcommand{\pending}{\textbf{\color{orange}[PENDING]}}
\newcommand{\fatal}{\textbf{\color{red}[OPEN]}}

\title{\textbf{General Theory for Nonlinear Convolution Equations}}
\author{}
\date{\today}

\begin{document}

\maketitle
\tableofcontents
\bigskip

This file contains the background material on abstract nonlinear convolution
equations and spectral support forcing.  The material is separated from the
specific $S^1$ sharp restriction programme, which is kept in
\texttt{convolution-equations.tex}.



% =======================================================================
\section{Setup and Notation}
\label{sec:setup}
% =======================================================================

We work throughout with \emph{entire functions of finite exponential type}.
An entire function $u:\C\to\C$ has \emph{exponential type} $\tau\geq 0$ if
$\abs{u(z)}\leq Ce^{\tau\abs{z}}$ for some $C>0$ and all $z\in\C$.
The \emph{Paley--Wiener--Schwartz theorem} identifies this class:
$u$ is entire of type $\leq\tau$ with $u|_{\R}\in L^2(\R)$ if and only if
$\wh u\in L^2$ is supported in $[-\tau,\tau]$.
We write $\PW_\tau$ for the corresponding Hilbert space.

Fourier conventions:
\[
  \wh u(\xi)=\int_\R u(x)e^{-ix\xi}\,dx,
  \qquad
  u(x)=\frac{1}{2\pi}\int_\R \wh u(\xi)e^{ix\xi}\,d\xi,
\]
so $\widehat{u\ast v}=\wh u\cdot\wh v$ and
$\widehat{u\cdot v}=\tfrac{1}{2\pi}\wh u\ast\wh v$.

Throughout, $k\geq 1$ is a fixed positive integer, $c\in\C\setminus\{0\}$
is a fixed constant, and $f$ is an entire function of exponential type
$\tau_f>0$ whose restriction to~$\R$ lies in $L^2(\R)$.

% =======================================================================
\section{The Compact Convolution Equation}
\label{sec:two-eqns}
% =======================================================================

\subsection{The equation $u=c\,(u^k\ast f)$: compact fixed-point structure}
\label{sec:reversed-eqn}

We now turn to the equation
\begin{equation}  \label{eq:main}
  u = c\,(u^k\ast f).
\end{equation}
In Fourier space:
\begin{equation}  \label{eq:main-fourier}
  \wh u(\xi) = c\,\wh u^{\ast k}(\xi)\cdot\wh f(\xi).
\end{equation}

\begin{proposition}[Type self-consistency]
  Any solution $u\in\PW_\tau$ of~\eqref{eq:main} satisfies $\tau\leq\tau_f$.
\end{proposition}

\begin{proof}
  The right side of~\eqref{eq:main-fourier} is supported in
  $\supp(\wh u^{\ast k})\cap\supp(\wh f)\subseteq[-\tau_f,\tau_f]$.
\end{proof}

This is the key contrast: the equation $u=c(u^k\ast f)$ does \emph{not}
force $\tau=0$; solutions of positive exponential type exist as long as
$\tau\leq\tau_f$.

\begin{proposition}[Compactness of $T$ on Paley--Wiener spaces]
  \label{prop:T-compact-PW}
  Define $T:\PW_{\tau_f}\to\PW_{\tau_f}$ by $T(v)=c(v^k\ast f)$.  Then
  $T$ maps bounded sets in $\PW_\tau$ (for any $\tau>\tau_f$) to precompact
  sets in $\PW_\tau$.  In particular, $T:\PW_\tau\to\PW_\tau$ is a compact
  operator.
\end{proposition}

\begin{proof}
  Since $v\in\PW_{\tau_f}$ implies $v^k\in\PW_{k\tau_f}$, the function
  $v^k\ast f$ has Fourier support in
  $\supp(\wh{v^k})\cap\supp(\wh f)\subseteq[-\tau_f,\tau_f]$.
  Hence $T$ maps $\PW_\tau$ into $\PW_{\tau_f}$, and the inclusion
  $\PW_{\tau_f}\hookrightarrow\PW_\tau$ is compact by the Rellich-type
  theorem for Paley--Wiener spaces.
\end{proof}

\subsection{Fredholm structure and local manifold theorem}

At a solution $u_0$ of~\eqref{eq:main}, the Fr\'echet derivative of
$G(u):=u-c(u^k\ast f)$ is
\[
  DG(u_0)[h] = h - ck(u_0^{k-1}h\ast f) =: (I-K)[h],
\]
where $K:h\mapsto ck(u_0^{k-1}h\ast f)$ is compact by the same argument.
By Riesz--Schauder, $I-K$ is \textbf{Fredholm of index zero}.

\begin{theorem}[Local manifold structure]
  \label{thm:local-mfd}
  Let $u_0$ be a solution of~\eqref{eq:main} at which $DG(u_0)=I-K$ is
  invertible.  Then $u_0$ is an isolated solution.  More generally, the
  solution set near $u_0$ is a smooth manifold of dimension
  $\dim\ker(I-K)$, which is finite.
\end{theorem}

% =======================================================================
\section{Spectral Support Forcing}
\label{sec:spectral-forcing}
% =======================================================================

We specialise to $f=\wh\sigma$, where $\sigma$ is the surface measure on
$S^{n-1}\subset\R^n$.  By the Paley--Wiener theorem, $f$ is entire of
exponential type~$1$, and $\wh f=(2\pi)^n\sigma$ as a distribution.

\begin{proposition}[Spectral support forcing]
  \label{prop:support-forcing}
  If $u$ is entire of exponential type satisfying $u=c(u^k\ast\wh\sigma)$,
  then $\wh u$ is a distribution supported on~$S^{n-1}$.
  The same conclusion holds for the equation $u=c(|u|^{2(k-1)}u\ast\wh\sigma)$.
\end{proposition}

\begin{proof}
  From~\eqref{eq:main-fourier}, $\wh u=c\,\wh u^{\ast k}\cdot(2\pi)^n\sigma$.
  The right side is supported on $\supp\sigma=S^{n-1}$.

  For the $|u|^{2(k-1)}u$ variant: writing $|u|^{2(k-1)}u=u^k\bar u^{k-1}$
  in Fourier space involves $(k-1)$-fold convolutions of $\wh{\bar u}$, where
  $\wh{\bar u}(\xi)=\overline{\wh u(-\xi)}$.  Since $S^{n-1}$ is symmetric
  under $\xi\mapsto-\xi$, the support of all convolved terms lies in the
  sumset of copies of $\pm S^{n-1}$, and multiplication by $\sigma$ restricts
  back to $S^{n-1}$.
\end{proof}

\begin{remark}
  In dimension $n\geq 2$, the conclusion $\supp\wh u\subseteq S^{n-1}$ means
  $u$ is a \emph{Herglotz wave function}, satisfying $(-\Delta-1)u=0$.
\end{remark}

\subsection{Reduction to a density equation on the sphere}

Write $\wh u=\rho\,d\sigma$ for $\rho\in L^2(S^{n-1})$, so
\[
  u(x) = \int_{S^{n-1}}\rho(\omega)e^{ix\cdot\omega}\,d\sigma(\omega).
\]
The equation $u=c(u^k\ast\wh\sigma)$ then reduces to the
\emph{spherical fixed-point equation}
\begin{equation}  \label{eq:sphere-eq}
  \rho(\omega) = c\int_{V_k(\omega)}\rho(\omega_1)\cdots\rho(\omega_k)
    \,d\sigma_{\mathrm{fiber}}(\omega_1,\ldots,\omega_k),
  \qquad \sigma\text{-a.e.}\ \omega\in S^{n-1},
\end{equation}
where $V_k(\omega)=\{(\omega_1,\ldots,\omega_k)\in(S^{n-1})^k:
\omega_1+\cdots+\omega_k=\omega\}$ is the incidence fibre.

\begin{warning}[Incidence fibre $\neq$ spherical convolution]
  \label{warn:not-convolution}
  The operator $T(\rho)(\omega)=\int_{V_k(\omega)}\rho^{\otimes k}\,
  d\sigma_\mathrm{fiber}$ is \emph{not} the same as convolution by the
  spherical measure.  In particular, the Funk--Hecke theorem does not
  diagonalise~$T$ in spherical harmonics.  Any reduction to finitely many
  harmonic coefficients requires separate justification.
\end{warning}

% =======================================================================
\section{The One-Dimensional Case: Complete Solution}
\label{sec:1D}
% =======================================================================

For $n=1$, $S^0=\{-1,+1\}$, $\sigma=\delta_{-1}+\delta_{+1}$, and
$f(x)=\wh\sigma(x)=2\cos x$.  Proposition~\ref{prop:support-forcing} forces
$\wh u=a\delta_{-1}+b\delta_{+1}$, so $u(x)=ae^{-ix}+be^{ix}$.

\subsection{Algebraic reduction}

Expanding $u^k$ and computing $u^k\ast f$ picks out frequencies
$\pm 1$, which exist only when $k$ is \textbf{odd}.

\begin{proposition}[Parity obstruction]
  If $k$ is even, the equation $u=c(u^k\ast f)$ with $f=2\cos$ has only
  $u\equiv 0$.
\end{proposition}

For $k$ odd, setting $\beta_k:=\binom{k}{(k-1)/2}$ and comparing coefficients
of $e^{\mp ix}$ gives the single constraint
\begin{equation}  \label{eq:constraint}
  ab = \lambda_k,
  \qquad
  \lambda_k := \left(\frac{1}{2\pi c\beta_k}\right)^{2/(k-1)}.
\end{equation}

\begin{theorem}[Complete solution in 1D]  \label{thm:main-1d}
  \proven\quad
  Let $k\geq 3$ be odd, $c\in\C\setminus\{0\}$, $f=2\cos$.
  The solution set of $u=c(u^k\ast f)$ in the Paley--Wiener class is
  \[
    \{0\}
    \;\cup\;
    \bigsqcup_{j=0}^{(k-3)/2}
      \{(a,b)\in\C^2:ab=\omega^j\lambda_k\},
    \qquad \omega=e^{4\pi i/(k-1)}.
  \]
  This is a finite union of $(k-1)/2$ smooth one-dimensional complex
  manifolds (each $\cong\C^*$), plus the origin.
  Each component is exactly the translation orbit of any single non-trivial
  solution on it: the translation $u(x)\mapsto u(x-t)$ acts by
  $(a,b)\mapsto(e^{-it}a,e^{it}b)$, preserving~$ab$.
\end{theorem}

\begin{example}[$k=3$]
  $\beta_3=3$.  Constraint: $ab=(6\pi c)^{-1}$.  One component, the
  translation orbit of $u(x)=(6\pi c)^{-1/2}\cos x$.
\end{example}

\begin{example}[$k=5$]
  $\beta_5=10$.  Constraint: $ab=(20\pi c)^{-1/2}$.  Two components
  (the scaling group $\mu_2=\{1,-1\}$ generates a second hyperbola).
\end{example}

\begin{remark}[Real solutions]
  For real-valued solutions $b=\bar a$, so $ab=|a|^2\geq 0$.  Only the
  hyperbola with $\lambda_k\in\R_{>0}$ contributes; the real solution
  manifold is the circle $|a|=|\lambda_k|^{1/2}$, parametrised by
  translations.
\end{remark}

% =======================================================================
% PART II: THE S^1 RESTRICTION PROBLEM
% =======================================================================


\end{document}
